\documentclass[11pt,a4paper]{article}
\usepackage[utf8]{inputenc}
\usepackage[T1]{fontenc}
\usepackage[french]{babel}
\usepackage[left=2cm,right=2cm,top=2cm,bottom=2cm]{geometry}
\usepackage{hyperref}
\usepackage{fontawesome5}

\pagestyle{empty}
\setlength{\parindent}{0pt}
\setlength{\parskip}{1em}

\begin{document}

% --- En-tête ---
\begin{center}
    {\Large \textbf{Loïc Aron MBASSI EWOLO}} \\
    \vspace{0.1cm}
    \small
    \faEnvelope\ \href{mailto:loicaron.mbassiewolo@ensea.fr}{loicaron.mbassiewolo@ensea.fr} \quad
    \faPhone\ +33 7 59 52 33 98 \quad
    \faGithub\ \href{https://github.com/Nameless0l}{Nameless0l}
\end{center}

\vspace{0.5cm}

\textbf{Objet :} Candidature Stage - Assistant Data \& AI Engineer (Copilot Agentique)

Madame, Monsieur,

Je candidate pour rejoindre le DataLab du Crédit Agricole sur le sujet des Copilotes Agentiques. Ce stage m'intéresse car il cible précisément mon domaine d'intérêt et est en lien avec mes anciens projets : l'outillage du développeur par l'IA.

Mon profil répond concrètement à vos missions pour trois raisons :

\textbf{1. Je sais construire des générateurs de code (Mission Prototype).}
Votre offre demande de générer des projets à partir de templates. J'ai déjà réalisé ce travail \textbf{disponible en opensource} avec mon projet \textit{\textbf{Laravel API Generator}}. J'ai dû résoudre les problèmes de parsing (UML/XML), d'injection de contexte dans les LLM et de respect de l'architecture cible. Je suis donc  immédiatement assez prêt sur ce volet technique.

\textbf{2. Je connais les architectures Agentiques (Mission Développement).}
Au-delà de la simple complétion, j'ai expérimenté l'approche Agentique (Google ADK, LangChain) en concevant un système multi-agents autonome. Je comprends les défis liés à la mémoire contextuelle et à l'exécution d'outils externes, essentiels pour dépasser les limites des copilotes statiques actuels.

\textbf{3. J'ai une rigueur scientifique pour l'évaluation (Mission Benchmark).}
Mon expérience d'Assistant de Recherche au \textbf{Mila (Institut Québécois d'IA) au Canada} sur la généralisation des réseaux de neurones (Grokking) m'a formé à l'analyse critique des modèles. J'y ai implémenté des pipelines de tests statistiques sous PyTorch et reproduit des expériences SOTA. Cette rigueur sera un atout majeur pour réaliser l'inventaire et le benchmark comparatif des solutions que vous attendez.

En parallèle, mon expérience freelance m'a formé aux standards industriels (CI/CD, Docker). Je ne ferai pas seulement du code de recherche, mais des outils intégrables dans vos pipelines de production.

Disponible dès avril 2026, je serais ravi de vous présenter mes projets lors d'un entretien.

Cordialement,

\vspace{0.5cm}

\textbf{Loïc Aron MBASSI EWOLO}

\end{document}